\begin{center}
\begin{tikzpicture}[x=.122\linewidth,y=6mm,every node/.style={font=\scriptsize,align=center}]
\fill[chFourTint] (0,5) rectangle (6.7,6);
\draw (0,0) rectangle (7.7,6);
\draw (2.8,0)--(2.8,6) (4.3,1)--(4.3,6) (5.4,1)--(5.4,6) (6.7,0)--(6.7,6);
\draw (0,1)--(6.7,1) (0,2)--(6.7,2) (0,5)--(6.7,5);
\draw (2.8,3)--(6.7,3) (2.8,4)--(6.7,4);
\node at (1.4,5.5) {Department};\node at (3.55,5.5) {Section};
\node at (4.85,5.5) {Male};\node at (6.05,5.5) {Female};
\node at (1.4,3.5) {H\textsubscript{2}O};\node at (1.4,1.5) {Business};\node at (1.4,.5) {Total};
\node at (3.55,4.5) {10};\node at (4.85,4.5) {10};\node at (6.05,4.5) {11};
\node at (3.55,3.5) {11};\node at (4.85,3.5) {12};\node at (6.05,3.5) {14};
\node at (3.55,2.5) {12};\node at (4.85,2.5) {15};\node at (6.05,2.5) {17};
\node at (3.55,1.5) {15};\node at (4.85,1.5) {16};\node at (6.05,1.5) {17};
\node at (4.75,.5) {420};\node at (7.2,3) {420};
\end{tikzpicture}
\captionof{figure}{অনুশীলন 21: প্রদত্ত কোড অনুযায়ী বিভাগভিত্তিক সংখ্যার ছক}
\end{center}
\textbf{চিত্র ব্যাখ্যা:} H\textsubscript{2}O ঘর তিন সারি জুড়ে আছে। ডান পাশের 420 ঘর শিরোনামসহ ছয় সারি এবং নিচের 420 ঘর তিন কলাম জুড়ে বিস্তৃত। Business-এর জন্য একটি তথ্যের সারি আছে। 420 এখানে প্রদত্ত নমুনা মান; পাশের সংখ্যাগুলোর গণিতসম্মত যোগফল নয়।
